% Reserved for future report appendices. Internal notes below are not typeset.
% Internal coverage and review notes (not typeset):
% \appendix
% \section{Working coverage and review notes}\label{sec:coverage}
% This appendix is a revision aid, not a question-by-question organization of the paper. ``Planned'' identifies evidence still to be collected. Q12 appears before the specification's Section 3 heading and is included in the core scope.
% \begin{center}
% \small
% \begin{tabularx}{\linewidth}{@{}l p{0.24\linewidth} X@{}}
% \toprule
% Question & Location & Current status \\
% \midrule
% Q1 & Section~\ref{sec:scaling} & Scaling derivation and PMF/CCDF distinction included. \\
% Q2 & Section~\ref{sec:rules} & SOC mechanisms explained; empirical evidence planned. \\
% Q3 & Section~\ref{sec:boundaries} & Finite/infinite cases discussed; infinite-case result cited. \\
% Q4 & Sections~\ref{sec:samplingtheory}, \ref{sec:design}, \ref{sec:results} & Naive prediction recorded; measured comparison planned. \\
% Q5 & Sections~\ref{sec:boundaries}, \ref{sec:design} & State concepts distinguished; diagnostics planned. \\
% Q6 & Sections~\ref{sec:samplingtheory}, \ref{sec:design} & Recurrence and sampling explained; comparison planned. \\
% Q7 & Sections~\ref{sec:observables}, \ref{sec:distributions}, \ref{sec:results} & Four definitions and distribution/fit plans included. \\
% Q8 & Sections~\ref{sec:correlations}, \ref{sec:results} & Predictions recorded; all six comparisons planned. \\
% Q9 & Section~\ref{sec:matrix} & Four conditions interpreted, with necessary caveats. \\
% Q10 & Section~\ref{sec:matrix} & Matrix, threshold, and boundary sums verified algebraically. \\
% Q11 & Section~\ref{sec:abelian} & Commutation proved for simultaneously unstable sites. \\
% Q12 & Section~\ref{sec:abelian} & Termination and uniqueness proved for the finite open grid. \\
% \bottomrule
% \end{tabularx}
% \end{center}
% 
% \paragraph{Report self-review.}
% \emph{Contribution:} Is the scope clear? Yes: a study of the baseline model, without a novelty claim.
% \emph{Clarity:} Are terms and narrative consistent? Model, theory, and proposed evidence are separated, with fixed notation and sample definitions.
% \emph{Experimental strength:} Are claims supported by measurements? None are claimed; every empirical conclusion remains pending.
% \emph{Evaluation completeness:} Are all core questions covered? Yes, at theory or planned-evidence level; figures, analyses, and an optional infinite-lattice proof remain TODOs.
% \emph{Method soundness:} Are important assumptions explicit? Boundary rules, legal schedules, finite support, and dependent sampling are specified; implementation and pilot decisions still need validation.
% 
% Claim-evidence map for revision (not experimental findings):
% Claim: Rescaling under the stated assumptions implies a power law.
% Evidence: algebra in Section 2.1. Status: supported under those assumptions.
% Claim: Finite-grid stabilization terminates and is unique.
% Evidence: potential and first-excess proofs in Section 3.2. Status: supported.
% Claim: Stationary recurrent heights are not independent uniform draws.
% Evidence: standard recurrent-state theory, Jarai Section 2. Status: supported.
% Claim: Limiting recurrent mean height is 17/8.
% Evidence: Levine and Peres. Status: literature result, not measured here.
% Claim: Finite-support inputs on the infinite grid stabilize.
% Evidence: Levine and Propp. Status: cited result; optional proof remains TODO.
% Claim: Particular exponents/correlations or computational speedups occur.
% Evidence: none. Status: not asserted; experiments required.
% TODO(final review): Resolve remaining references and proof notes, replace
% empirical TODOs, inspect all figures, and remove working review notes only
% after core-question coverage and the 20-page limit have been checked again.
